\chapter*{INTRODUCTION}\addcontentsline{toc}{chapter}{INTRODUCTION}\label{introduction}

\section{CONTEXTE GÉNÉRAL DE L'ÉTUDE}
\label{contexte-general-de-letude}

L'accès à une eau potable sûre et disponible de manière régulière demeure
un enjeu important dans de nombreux pays en développement. Selon le
Programme commun de suivi de l'Organisation mondiale de la Santé (OMS)
et du Fonds des Nations Unies pour l'enfance (UNICEF), près de
2,2 milliards de personnes ne disposaient pas, en 2022, d'un service
d'eau potable géré en toute sécurité \cite{ref1}. Cette situation montre
que l'amélioration de l'accès à l'eau potable ne dépend pas uniquement
de la disponibilité de la ressource, mais également de la capacité à
assurer une gestion efficace et durable des infrastructures de
distribution.

En République Démocratique du Congo, cette problématique demeure
particulièrement importante dans les centres urbains où l'augmentation
de la population et l'extension des zones d'habitation exercent une
pression croissante sur les infrastructures existantes. Dans ce
contexte, la continuité de la desserte dépend notamment de l'état des
conduites, du fonctionnement des équipements hydrauliques, de la
disponibilité de l'eau et de la capacité des services techniques à
identifier et à traiter rapidement les anomalies. Les difficultés de
surveillance et de maintenance peuvent ainsi se traduire par des
interruptions de service, des pertes d'eau et des interventions
principalement effectuées après l'apparition d'une défaillance.

La ville de Bukavu présente ces enjeux dans un contexte urbain
particulier, caractérisé par une forte densité de population et une
extension progressive des zones d'habitation. La commune de Kadutu,
qui constitue l'une des communes de la ville, est directement concernée
par ces difficultés. Les observations réalisées au cours de la phase
exploratoire de cette étude, associées aux informations disponibles
dans les sources documentaires, montrent notamment l'existence de
perturbations de la desserte en eau potable dans certains secteurs de
Kadutu. En juillet 2022, l'Agence Congolaise de Presse rapportait par
exemple que certaines avenues de la commune rencontraient des
difficultés d'approvisionnement en eau distribuée par la REGIDESO,
avec une situation particulièrement préoccupante durant la saison
sèche \cite{ref127}.

Dans ce contexte, la gestion efficace du réseau de fourniture d'eau
potable de Kadutu nécessite non seulement des infrastructures physiques
adaptées, mais également des moyens permettant de mieux connaître
l'état du réseau et de suivre son fonctionnement. Or, les informations
relatives aux différents éléments du réseau, à leur état, à leur
localisation et à leur fonctionnement ne sont pas toujours regroupées
dans un environnement numérique permettant une exploitation et une
analyse centralisées. Cette situation rend plus difficile la
visualisation globale du réseau et l'identification rapide des secteurs
susceptibles de nécessiter une intervention.

Les progrès récents de l'intelligence artificielle et des technologies
de représentation numérique offrent cependant de nouvelles possibilités
pour répondre à ce type de difficulté. Le jumeau numérique permet de
représenter virtuellement un système physique et de faire évoluer cette
représentation à partir des données disponibles. Associé à des
techniques d'intelligence artificielle, il peut être utilisé pour
analyser les données du réseau, identifier certaines anomalies et
simuler des situations de fonctionnement. De telles possibilités sont
particulièrement intéressantes dans le contexte de Kadutu, où
l'amélioration de la connaissance du réseau et de sa surveillance peut
constituer un appui à la maintenance et à la prise de décision
\cite{ref9}.

C'est donc dans le contexte spécifique du réseau de fourniture d'eau
potable de la commune de Kadutu que s'inscrit la présente étude. Elle
vise à étudier la possibilité de mettre en place un jumeau numérique
intégrant des techniques d'intelligence artificielle afin de disposer
d'une représentation numérique du réseau, de faciliter sa surveillance,
d'identifier certaines anomalies et de fournir un support à la
planification des interventions de maintenance.


\section{PROBLÉMATIQUE}
\label{problematique}

Dans la commune de Kadutu, la gestion du réseau de fourniture d'eau
potable est confrontée à plusieurs difficultés qui limitent la capacité
à assurer un suivi efficace de son fonctionnement. Les perturbations
observées dans la desserte constituent l'un des principaux problèmes.
Certaines zones peuvent connaître des interruptions ou une distribution
irrégulière de l'eau, notamment dans certaines périodes de l'année.
Ces situations nécessitent de pouvoir déterminer rapidement les parties
du réseau concernées et d'identifier les facteurs susceptibles
d'expliquer la perturbation.

Un second problème concerne la connaissance et la surveillance du
réseau. Un réseau de fourniture d'eau est constitué de plusieurs
éléments interdépendants, notamment des conduites, des nœuds, des
réservoirs, des vannes et des points de desserte. Lorsque les
informations relatives à ces éléments ne sont pas représentées dans un
environnement numérique centralisé et facilement exploitable, il
devient plus difficile d'obtenir une vision globale de la structure du
réseau et de suivre son évolution. Cette insuffisance peut compliquer
l'analyse d'une anomalie et la localisation de la zone nécessitant une
intervention.

À cette difficulté s'ajoute la détection des anomalies pouvant affecter
le fonctionnement du réseau. Une baisse anormale de pression, une
variation inhabituelle du débit, une rupture de conduite ou une fuite
peuvent perturber la desserte et entraîner des pertes d'eau. Lorsque
ces phénomènes ne sont pas détectés suffisamment tôt, leur localisation
et leur prise en charge peuvent devenir plus difficiles. La disponibilité
de données structurées et leur analyse automatique pourraient donc
permettre d'améliorer la détection des situations anormales.

La maintenance constitue également un enjeu important. En l'absence
d'un outil permettant de centraliser les informations sur le réseau et
d'analyser son état, les interventions peuvent être davantage orientées
vers la réaction à une panne déjà apparue que vers son anticipation.
Cette situation limite la possibilité de hiérarchiser les zones
nécessitant une attention particulière et de préparer les interventions
sur la base d'informations relatives au fonctionnement du réseau.

Ainsi, le problème central identifié dans le cadre de cette étude n'est
pas uniquement l'existence de perturbations dans la fourniture d'eau,
mais également la difficulté à disposer d'un outil numérique permettant
de représenter le réseau de Kadutu, de suivre son fonctionnement,
d'identifier les anomalies et d'appuyer les décisions de maintenance.
L'insuffisance d'une telle représentation dynamique constitue une limite
pour une gestion fondée sur les données.

Le projet de jumeau numérique proposé dans cette étude cherche donc à
répondre à ces insuffisances en développant une représentation
numérique du réseau de fourniture d'eau potable de Kadutu. L'intégration
de techniques d'intelligence artificielle doit permettre d'exploiter les
données disponibles afin de contribuer à la détection des anomalies et
à l'analyse du fonctionnement du réseau. Le système envisagé doit
également permettre de visualiser les différents composants du réseau,
de simuler certaines situations et de fournir des informations utiles
à l'identification des secteurs nécessitant une intervention.

Dès lors, la question fondamentale qui guide cette étude est la
suivante :

\textbf{Comment la conception et le développement d'un jumeau numérique
	intégrant des techniques d'intelligence artificielle peuvent-ils
	contribuer à améliorer la surveillance, la détection des anomalies et
	la maintenance du réseau de fourniture d'eau potable de la commune de
	Kadutu ?}

\section{HYPOTHÈSES DE RECHERCHE}
\label{hypotheses-de-recherche}

\subsection{Hypothèse générale}
\label{hypothese-generale}

La mise en œuvre d'un \textbf{jumeau numérique intégrant des techniques
	d'intelligence artificielle} devrait améliorer la surveillance du
réseau de fourniture d'eau potable de la commune de \textbf{Kadutu},
faciliter la détection des anomalies et contribuer à une maintenance
plus efficace grâce à l'exploitation des données du réseau.

\subsection{Hypothèses spécifiques}
\label{hypotheses-specifiques}

Afin de répondre à la problématique de recherche, les hypothèses
spécifiques suivantes sont formulées :

\subsection{Première hypothèse spécifique}
\label{premiere-hypothese-specifique}

La modélisation numérique du réseau de fourniture d'eau potable de
\textbf{Kadutu} permettrait d'améliorer sa représentation, sa
visualisation et le suivi de ses différents composants.

\subsection{Deuxième hypothèse spécifique}
\label{deuxieme-hypothese-specifique}

L'intégration des données relatives au fonctionnement du réseau
permettrait d'améliorer la surveillance de celui-ci et de faciliter
l'identification des variations anormales de ses paramètres.

\subsection{Troisième hypothèse spécifique}
\label{troisieme-hypothese-specifique}

L'utilisation de techniques d'\textbf{intelligence artificielle}
permettrait de détecter automatiquement certaines anomalies du réseau
et de contribuer à l'anticipation des défaillances.

\subsection{Quatrième hypothèse spécifique}
\label{quatrieme-hypothese-specifique}

L'intégration du jumeau numérique et de l'intelligence artificielle
pourrait fournir aux gestionnaires des informations utiles à la
planification et à la priorisation des interventions de maintenance.

\section{OBJECTIFS DE L'ÉTUDE}
\label{objectifs-de-letude}

\subsection{Objectif général}
\label{objectif-general}

L'objectif général de cette étude est de \textbf{concevoir et développer
	un jumeau numérique intégrant des techniques d'intelligence artificielle}
pour améliorer la surveillance, la détection des anomalies et la
maintenance du réseau de fourniture d'eau potable de la commune de
Kadutu.

\subsection{Objectifs spécifiques}
\label{objectifs-specifiques}

\noindent{Pour atteindre cet objectif général, la présente étude vise à :}

\noindent\textbf{Analyser le réseau de fourniture d'eau potable de la
	commune de Kadutu} afin d'identifier ses principales composantes, son
fonctionnement et les difficultés liées à sa surveillance et à sa
maintenance ;

\noindent\textbf{Modéliser numériquement le réseau de fourniture d'eau
	potable de la commune de Kadutu} en représentant ses principaux
composants et leurs relations ;

\noindent\textbf{Mettre en place un mécanisme de collecte et
	d'intégration des données du réseau} afin de faciliter le suivi de son
état et de ses paramètres de fonctionnement ;

\noindent\textbf{Intégrer des techniques d'intelligence artificielle}
pour analyser les données du réseau et détecter certaines anomalies
susceptibles d'affecter son fonctionnement ;

\noindent\textbf{Développer une plateforme de jumeau numérique}
permettant de visualiser le réseau de la commune de Kadutu, de suivre
son état, de simuler certaines situations et de fournir des
informations utiles à la planification des interventions de
maintenance.

\section{INTÉRÊT DE
L\textquotesingle ÉTUDE}\label{intuxe9ruxeat-de-luxe9tude}

\subsection{Intérêt
scientifique}\label{intuxe9ruxeat-scientifique}

Sur le plan scientifique, cette étude contribue à enrichir les travaux
portant sur l\textquotesingle application des technologies du
\textbf{jumeau numérique}, de l\textquotesingle Internet des objets et
de l\textquotesingle intelligence artificielle dans le domaine de la
gestion des réseaux de distribution d\textquotesingle eau potable. Elle
participe également à l\textquotesingle adaptation de ces technologies
au contexte des pays en développement, et plus particulièrement à celui
de la République Démocratique du Congo, où peu de recherches ont
jusqu\textquotesingle à présent porté sur leur mise en œuvre dans les
infrastructures hydrauliques urbaines.

\subsection{Intérêt pratique}\label{intuxe9ruxeat-pratique}

Sur le plan pratique, cette étude vise à proposer une solution
susceptible d\textquotesingle améliorer la surveillance, la maintenance
et la planification du réseau de fourniture d\textquotesingle eau
potable de la commune de Kadutu. Les résultats attendus pourront servir
d\textquotesingle outil d\textquotesingle aide à la décision pour la
REGIDESO et les autorités compétentes, en facilitant la détection des
anomalies, la réduction des pertes d\textquotesingle eau,
l\textquotesingle optimisation des interventions techniques et
l\textquotesingle amélioration de la qualité du service rendu aux
usagers.

\section{DÉLIMITATION DE L'ÉTUDE}
\label{delimitation-de-letude}

\subsection{Délimitation spatiale}
\label{delimitation-spatiale}

La présente étude est réalisée dans la \textbf{commune de Kadutu},
située dans la ville de Bukavu, province du Sud-Kivu, en République
Démocratique du Congo. Elle porte spécifiquement sur le réseau de
fourniture d'eau potable de cette commune et vise à y expérimenter une
approche basée sur le jumeau numérique.

\subsection{Délimitation temporelle}
\label{delimitation-temporelle}

Cette recherche couvre la période allant de \textbf{2025 à 2026},
correspondant aux différentes étapes de l'étude, notamment la collecte
et l'analyse des données relatives à la structure, à la localisation et
aux caractéristiques techniques du réseau de fourniture d'eau potable
de la commune de Kadutu, ainsi qu'à la conception, au développement et à
l'évaluation du jumeau numérique proposé.

\subsection{Délimitation thématique}
\label{delimitation-thematique}

La présente étude s'inscrit dans le domaine de l'informatique appliquée,
plus précisément dans les thématiques du \textbf{jumeau numérique} et
de l'\textbf{intelligence artificielle} appliqués à la surveillance,
à la détection des anomalies et à la maintenance du réseau de
fourniture d'eau potable de la commune de Kadutu.

\subsection{Délimitation méthodologique}
\label{delimitation-methodologique}

Sur le plan méthodologique, l'étude porte sur l'analyse du réseau
existant, la collecte et la structuration des données disponibles, la
modélisation numérique de ses principaux composants, puis la conception
et le développement d'un prototype de jumeau numérique intégrant des
techniques d'intelligence artificielle. L'évaluation porte sur les
fonctionnalités de représentation du réseau, de surveillance, de
simulation et de détection des anomalies. Les résultats obtenus sont
évalués dans le cadre du prototype développé et ne prétendent pas
constituer une généralisation à l'ensemble des réseaux de fourniture
d'eau potable de la ville de Bukavu.

\section{SUBDIVISION DU MÉMOIRE}\label{subdivision-du-muxe9moire}

Le présent mémoire est subdivisé en \textbf{quatre chapitres} :

\begin{itemize}
\item
  \textbf{Chapitre I :} Revue de la littérature et fondements
  théoriques.
\item 
  \textbf{Chapitre II :} Matériels et méthodes.
\item
  \textbf{Chapitre III :} Conception et développement du jumeau
  numérique intégrant l\textquotesingle intelligence artificielle.
\item
  \textbf{Chapitre IV :} Présentation, analyse et discussion des
  résultats.
\end{itemize}

\clearpage
% ------- Fin de l'introduction